\hypertarget{python-ux6559ux7a0b}{%
\subsection{Python 教程}\label{python-ux6559ux7a0b}}

这是小白的 Python 新手教程，具有如下特点：

Python
是一种计算机程序设计语言。你可能已经听说过很多种流行的编程语言，比如非常难学的
C 语言，非常流行的 Java 语言，适合初学者的 Basic 语言，适合网页编程的
JavaScript 语言等等。

那 Python 是一种什么语言？

首先，我们普及一下编程语言的基础知识。用任何编程语言来开发程序，都是为了让计算机干活，比如下载一个
MP3，编写一个文档等等，而计算机干活的 CPU
只认识机器指令，所以，尽管不同的编程语言差异极大，最后都得 ``翻译'' 成
CPU
可以执行的机器指令。而不同的编程语言，干同一个活，编写的代码量，差距也很大。

比如，完成同一个任务，C 语言要写 1000 行代码，Java 只需要写 100 行，而
Python 可能只要 20 行。

所以 Python 是一种相当高级的语言。

你也许会问，代码少还不好？代码少的代价是运行速度慢，C 程序运行 1
秒钟，Java 程序可能需要 2 秒，而 Python 程序可能就需要 10 秒。

那是不是越低级的程序越难学，越高级的程序越简单？表面上来说，是的，但是，在非常高的抽象计算中，高级的
Python 程序设计也是非常难学的，所以，高级程序语言不等于简单。

但是，对于初学者和完成普通任务，Python 语言是非常简单易用的。连 Google
都在大规模使用 Python，你就不用担心学了会没用。

用 Python 可以做什么？可以做日常任务，比如自动备份你的
MP3；可以做网站，很多著名的网站包括 YouTube 就是 Python
写的；可以做网络游戏的后台，很多在线游戏的后台都是 Python
开发的。总之就是能干很多很多事啦。

Python 当然也有不能干的事情，比如写操作系统，这个只能用 C
语言写；写手机应用，只能用 Swift/Objective-C（针对 iPhone）和 Java（针对
Android）；写 3D 游戏，最好用 C 或 C++。

如果你是小白用户，满足以下条件：

\begin{itemize}
\item
  会使用电脑，但从来没写过程序；
\item
  还记得初中数学学的方程式和一点点代数知识；
\item
  想从编程小白变成专业的软件架构师；
\item
  每天能抽出半个小时学习。
\end{itemize}

不要再犹豫了，这个教程就是为你准备的！

准备好了吗？

 
 \begin{figure}[htp]
	\centering
	\includegraphics[width=0.6\linewidth]{fig/9229153429258240.png}
\end{figure}


\hypertarget{ux5173ux4e8eux4f5cux8005}{%
\subsubsection{关于作者}\label{ux5173ux4e8eux4f5cux8005}}

\href{http://weibo.com/liaoxuefeng}{廖雪峰}，十年软件开发经验，业余产品经理，精通
Java/Python/Ruby/Scheme/Objective C
等，对开源框架有深入研究，著有《Spring 2.0
核心技术与最佳实践》一书，多个业余开源项目托管在
\href{https://github.com/michaelliao}{GitHub}，欢迎微博交流：

使用窄屏手机的童鞋，请点击左上角 ``目录'' 查看教程：

 
 \begin{figure}[htp]
	\centering
	\includegraphics[width=0.6\linewidth]{fig/1311543585144897l.png}
\end{figure}


